\chapter*{ADVERTISEMENT.}
\thispagestyle{empty}
\addcontentsline{toc}{chapter}{ADVERTISEMENT}

\textsc{The} following pages were written by the desire of the London Consolidated
Society of Bookbinders, as expressed in the following resolution;---

\begin{quote}
\textit{Lodge Meeting, Nov.\ 7, 1859, Mr.\ T. Marston in the Chair.}
\\[1em]
\textsc{Resolved,}---That the Secretary of this Society be requested to write
an essay on the Philosophy of Trades' Unions, in reply to the argument against
them, to be published in the form of a pamphlet, explaining their true object
and intention.

(Moved by Mr.\ \textsc{W. Franklin.} Seconded by Mr.\ \textsc{Robt. Saunders.}
\end{quote}

The writer is aware of many imperfections; and doubtless many more, of which
he has not been aware, will be discovered by those who do him the honour of
perusal. Having been brought up to the trade, and for many years a workman,
he has stated what he knows to be true; and if it be said that his is only the
working man's view of the subject, it may be replied that such view is, of all
others, the one most needed; as the ``press'' has teemed ``with little else
than the employers' view of it; and who can form a true judgment when only one
side of a question is presented to him?

Although not required to do so, he has, previous to their going to the printer,
read the whole of the following pages to the Committee of his Society, who, in
the following resolution, expressed their opinion of them:---

\begin{quote}
\begin{flushright}
Committee Room, January 25th, 1860.
\end{flushright}

\textsc{Resolved,}---That the concluding part of the Essay on the Philosophy of
Trades' Unions and Strikes, now read, and those read at previous meetings of
this Committee, are in its opinion a faithful representation of their views on
the subject, and, as they believe, those of their trade generally.
\end{quote}

\begin{center}
\rule{10cm}{0.4pt}\\
\end{center}

What is written, it is, perhaps, not too much to say, may therefore be taken
to set forth, as far as they and those of their trade are to be regarded as a
sample, not simply their views, or those of the writer, but those of the
working classes generally.

\scalebox{1.5}{\ding{43}}~The members of our Society will doubtless be surprised at these pages being
published at their own office, after two or three first-class publishers being
named at the meeting which passed the above resolution. But amid the outcry for
free trade in labour there does not appear to be free trade in publishing.
Although excellent reasons were given for declining to publish this work, and
with the greatest civility, we cannot conceal from ourselves that such a work
as this is altogether ``tabooed'' by those who are termed ``respectable''
publishers; not that they would have any objection; but as it appears to us,
they fear the ``public opinion'' of the class to which they belong, which we
need not say is dead against Trades' Unions. We hope we are mistaken, yet we
cannot but think that, however essential it may be to know both sides of a
question, that on this subject the middle class do not wish that any but their
own side should be known.

We, however, should ill deserve, as ``English workmen,'' the high character
given of us by the ``Quarterly'' reviewer, if we felt at a loss from this
circumstance. We therefore publish the work ourselves. But as we are aware
that we cannot do it so effectively as ``the trade,'' we beg the co-operation
of those who feel interested in our success. Should there be any difficulty in
obtaining this work, twelve postage stamps sent to the publisher will insure
its transmission, post free, to any part of the United Kingdom.

\begin{flushleft}
\textit{No.\ 5, Racquet Court, Fleet-street, E.C.}
\end{flushleft}
